% Chapter 5: Measures and statistical analysis

\chapter{Measures and Statistical Analysis}
\label{ch:framework}

This chapter defines the outcomes, denominators, and uncertainty procedures used
to answer the two research questions.

\section{Analysis principles}
\label{sec:framework-principles}

We report the available sample, point estimate, and 95\% confidence interval
from a bootstrap clustered by source frame where applicable. The analysis plan, including its
comparisons, was fixed before the complete results were inspected. Every planned comparison is reported.

Differences and intervals were computed from unrounded values; the components
displayed in tables are rounded for presentation, so subtracting two printed
numbers can differ from the printed difference in the last digit.

Missing rationale text remained missing. The primary Stage~2 analysis uses a
paired record only when the tolerant parser recovered a non-empty final rationale span
from both source and intervention outputs. We report parser coverage before the
scientific estimates so that a change in denominator remains visible.

\section{RQ1: answer behaviour}
\label{sec:framework-rq1}

\subsection{Accuracy and gaps relative to source}
\label{sec:framework-accuracy}

Accuracy is the proportion of records whose Stage~1 prediction equals the gold
answer. With four choices, chance performance is 0.25. The five intervention
conditions are grey, mismatch, mask, noise, and mirror. For any one of these
conditions, the gap relative to source is
\[
  G_{\mathrm{condition}}
  =A_{\mathrm{source}}-A_{\mathrm{condition}},
\]
where $A_{\mathrm{source}}$ is accuracy with the source image and
$A_{\mathrm{condition}}$ is accuracy under the named intervention condition.
Positive values mean that accuracy was lower after the
image intervention. Both component accuracies accompany the gap because the
same difference can arise at different performance levels.

The \emph{Sufficiency Gap} is the gap for grey replacement:
\[
  \mathrm{Sufficiency\ Gap}
  =A_{\mathrm{source}}-A_{\mathrm{grey}}.
\]
It quantifies how much answer accuracy is lost when scene pixels are replaced
while the question, choices, checkpoint, and optional description remain. Grey
also changes the size and shape of the input (Section~\ref{sec:setup-conditions}),
so the gap measures this specific 256 by 256 replacement. The same grey
constructor was applied to every arm, but interactions with its size, shape, and
processed visual tokens may still differ between arms. A
smaller gap means that more information relevant to the answer
remains available under grey. In a description arm, the retained text
can provide that information, so a smaller gap is interpreted as behavioural
robustness.

The planned comparisons across arms compare \texttt{plain\_desc} with
\texttt{plain} and \texttt{point\_desc} with \texttt{point}. They subtract
the Sufficiency Gap without a description from the corresponding gap with a description.
Negative values mean that the description arm lost less accuracy.

\subsection{Correctness transitions}
\label{sec:framework-transitions}

Net accuracy can hide opposing movements for individual records. We therefore
classify each source and condition pair as correct to correct (C$\rightarrow$C),
correct to wrong (C$\rightarrow$W), wrong to correct (W$\rightarrow$C), or
wrong to wrong (W$\rightarrow$W). C$\rightarrow$W is reported relative to the
number correct under source, while W$\rightarrow$C is reported relative to the
number wrong under source and to all records.

Churn is C$\rightarrow$W plus W$\rightarrow$C. It measures how many records
changed correctness state even when the two directions cancel in aggregate
accuracy. Net movement is W$\rightarrow$C minus C$\rightarrow$W. The
accounting identity
\[
  A_{\mathrm{source}}-A_{\mathrm{condition}}
  =\frac{\left(\mathrm{C}\rightarrow\mathrm{W}\right)
         -\left(\mathrm{W}\rightarrow\mathrm{C}\right)}{N}
\]
was checked in every combination of arm and condition.

The qualitative W$\rightarrow$C audit uses a fixed sample of 16 cases, each
representing one arm under one condition. Twelve represent the conditions other than mismatch and four represent
mismatch under the shared map. Their identifiers were fixed before questions,
choices, or images were inspected. The audit is illustrative and does not estimate category
prevalence.

\section{RQ2: BERTScore Drift}
\label{sec:framework-rq2}

\subsection{Why compare paired rationales}
\label{sec:framework-drift-motivation}

RQ2 asks whether generated rationale text responds when visual evidence changes.
Reference similarity to a human rationale does not answer this question because
both source and intervention outputs could be equally plausible. We compare the
two model outputs from the same record, arm, checkpoint, question, choices, gold
answer, and decoding configuration.

Let $R_{i,\mathrm{source}}$ be the final rationale from the tolerant parser for
record $i$ under the source image, and let $R_{i,\mathrm{condition}}$ be the
final rationale for the same record under the named intervention condition. The
primary measure is
\[
  D_{i,\mathrm{condition}}
  =1-\mathrm{BERTScore}_{\mathrm{F1}}
  \left(R_{i,\mathrm{source}},R_{i,\mathrm{condition}}\right).
\]
We call $D_{i,\mathrm{condition}}$ \emph{BERTScore Drift}, or Drift after first
use. We convert BERTScore F1 similarity into a dissimilarity score by
subtracting the similarity value from 1. Larger values therefore mean greater
text change on this configured scale. Identical strings have Drift zero by
definition. Drift measures change in the generated text; truth, semantic
correctness, and hidden causal reasoning are separate questions.

\subsection{BERTScore configuration}
\label{sec:framework-bertscore-config}

We computed BERTScore F1 \citep{DBLP:conf/iclr/ZhangKWWA20}. Our implementation
used \texttt{bert\_score==0.3.12} with \textsc{roberta-large}, model snapshot
\texttt{722cf37b1afa9454edce342e7895e588b6ff1d59}, and layer~17.
Inverse document frequency weighting was
disabled, and baseline rescaling was disabled, so every reported value is raw.
Precision, recall, F1, and raw Drift were stored for every eligible pair
together with hashes of both text spans.

For each arm and intervention, we report the number of usable pairs, the mean,
and a 95\% confidence interval for the mean. Whether the ordering of the
conditions depends on scoring the final span is checked in
Appendix~\ref{sec:app-rq2-sensitivity}, which repeats the same pairs on the
reasoning span and the complete generation.

\section{Drift reference points}
\label{sec:framework-anchors}

Raw Drift values are easier to interpret when they are placed beside three
observed reference points.

\paragraph{Mirror.}
For each record, we compare the rationale generated with the source image with
the rationale generated after horizontally mirroring that image. The gold answer
remains fixed. Mirror preserves the scene content while reversing its
orientation, so it gives a reference point for a visual change that does not
remove or replace the scene.

\paragraph{Wrong-answer control.}
For the same fixed 200 records in each arm, we use the source image in both
answer conditions:

\begin{itemize}
  \item First, Stage~2 receives the source image and the gold answer and
        produces the source rationale.
  \item Second, Stage~2 receives the same source image and a predefined answer
        that differs from the gold answer and produces another rationale.
  \item We compute the Drift between these two rationales.
\end{itemize}

For the matched grey value, we compare the same source rationale with the
rationale generated from the grey image while the answer remains gold. The
paired comparison is
\[
  \Delta_{\mathrm{answer}}
  =D_{\mathrm{wrong\ answer}}-D_{\mathrm{matched\ grey}}.
\]
A positive value means that changing the supplied answer produced more rationale
change than grey replacement on the same records. The analysis on the reasoning
span, and the analysis on the final span with answer choices masked, check
whether this direction depends on direct lexical echo.

\paragraph{Reference using rationales from different records.}
Within each arm, we form a fixed one-to-one pairing of source-condition
rationales:

\begin{itemize}
  \item Take the source-condition rationale for each record.
  \item Pair it with the source-condition rationale from a different record in
        the same arm.
  \item Require the two records to come from different source frames.
  \item Compute Drift for every pair and average the resulting values.
\end{itemize}

The mean shows the scale of Drift when the two rationales describe different
records. We use it as an empirical reference point for this dataset and metric
configuration.

The one-to-one pairing used seed \texttt{20260829} for \texttt{plain},
\texttt{20260830} for \texttt{point}, \texttt{20260831} for
\texttt{plain\_desc}, and \texttt{20260832} for \texttt{point\_desc}. The
accepted permutation was the third draw for \texttt{plain}, the second for
\texttt{point} and \texttt{plain\_desc}, and the first for
\texttt{point\_desc}.

\section{Uncertainty and planned comparisons}
\label{sec:framework-statistics}

Several VCR questions can use the same movie frame, so records from one frame
are related. We therefore estimated uncertainty with a bootstrap clustered by
source frame \citep{DBLP:books/sp/EfronT93}.

\begin{itemize}
  \item \textbf{Sampling unit.} One source frame is one bootstrap unit, together
        with every analysed record that uses that frame. A full-population cell
        contains 2,400 frame units. Cells affected by a missing rationale contain
        2,398 or 2,399, and the 200-record wrong-answer control contains 185.

  \item \textbf{Resampling.} For each replicate, sample with replacement the
        same number of frame units as the analysed set contains. When a frame is
        selected more than once, all records linked to it receive the same
        multiplicity.

  \item \textbf{Recalculation.} Recompute the requested accuracy, transition
        rate, mean Drift, or difference between paired terms on the resampled
        data. A paired comparison uses the same sampled frame multiplicities for
        both terms.

  \item \textbf{Repetition.} Repeat the process 10,000 times with seed~42.

  \item \textbf{Confidence interval.} Use the 2.5th and 97.5th percentiles of
        the 10,000 replicate values as the endpoints of the 95\% interval.
\end{itemize}

An interval that includes zero leaves the direction of the difference unresolved
under this procedure. Establishing equivalence would require a dedicated
equivalence test.

The confidence intervals are unadjusted.
Section~\ref{sec:framework-principles} records that the comparison list was
fixed in advance, and every planned comparison is reported.

\begin{itemize}
  \item \textbf{RQ1.} Compare \texttt{plain\_desc} with \texttt{plain} and
        \texttt{point\_desc} with \texttt{point} for the Sufficiency Gap and for
        the grey C$\rightarrow$W rate.

  \item \textbf{RQ2 intervention values.} Report Drift for grey, mismatch, mask,
        noise, and mirror relative to source in every arm.

  \item \textbf{RQ2 planned pairwise comparisons.} Within each arm, compare grey
        with mirror, mismatch with mirror, grey with mismatch, and mask with
        noise. Across matched arms under grey, compare \texttt{plain\_desc} with
        \texttt{plain} and \texttt{point\_desc} with \texttt{point}. These are
        the comparisons reported in
        Appendix~\ref{sec:app-rq2-contrasts}.

  \item \textbf{Wrong-answer control.} Compare wrong-answer Drift with matched
        grey Drift on the fixed 200-record, 185-frame sample.
\end{itemize}

\section{PRED-A and Cross-Modal Consistency}
\label{sec:framework-pred-a}

The secondary predicted-answer regime (PRED-A) is measured by Cross-Modal
Consistency (CMC), which asks whether the answer Stage~2 was given can be
recovered from the rationale it wrote. We measure this by letting a text-only
model act as the judge. The judge receives the question, the four answer
choices, and the eligible final rationale span. These are its only inputs. It
selects one of the four choice indices, and recovery succeeds when that choice
matches the answer supplied to Stage~2.

The judge is \texttt{Qwen/Qwen2.5-3B-Instruct}, pinned to revision
\texttt{aa8e72537993ba99e69dfaafa59ed015b17504d1}. Inputs were truncated at 2,048
tokens and
scored in batches of 16. The reported CMC values are properties of this judge as
much as of the rationales, and a different judge would give different numbers.

% Keep the sentence that names R, U and N on the same page as the equations it
% introduces; a break here left the closing word alone above the display.
\begin{samepage}
Two rates follow from that, and they answer different questions. Writing $R$ for
the number of recoveries, $U$ for the usable final rationales, and $N$ for all
records,
\[
  \mathrm{CMC}_{\mathrm{conditional}} = \frac{R}{U},
  \qquad
  \mathrm{CMC}_{\mathrm{pipeline}} = \frac{R}{N}.
\]
\end{samepage}
Conditional CMC describes only the records for which a usable rationale exists.
Pipeline CMC divides by every record, so it carries generation coverage with it.
Where coverage is low the two diverge sharply, and reporting only the first
would describe the successful generations rather than the regime.

CMC records coherence between a supplied answer and the rationale written for it.
It does not measure visual grounding, answer correctness, rationale truth, or
faithfulness. Uniform random selection among four choices succeeds with
probability 0.25. This is the mathematical chance level of the four-way
decision. Because the judge also sees the question and answer choices, CMC may partly
reflect information available without the rationale; estimating that
contribution would require a question-and-choices-only baseline.
Appendix~\ref{sec:app-prompts} reproduces the judge prompt and states its
decoding rule.
